% ECONOMICS_PROBLEM: {"id": "F35-575", "title": "Aid that builds lasting capacity", "jel_code": "F35", "source_id": "OP-0375"}
% !TeX program = xelatex
\documentclass[11pt,a4paper]{article}
\usepackage[margin=23mm]{geometry}
\usepackage{fontspec}
\setmainfont{Latin Modern Roman}
\usepackage{amsmath,amssymb,mathtools}
\usepackage{xcolor,enumitem,booktabs,longtable,array,xurl,hyperref}
\definecolor{muted}{HTML}{53636C}
\hypersetup{colorlinks=true,urlcolor=blue,linkcolor=blue}
\setlength{\parindent}{0pt}
\setlength{\parskip}{5pt}
\newcommand{\R}{\mathbb R}
\newcommand{\E}{\mathbb E}
\newcommand{\N}{\mathbb N}
\newcommand{\ind}{\mathbf 1}
\newcommand{\Var}{\operatorname{Var}}
\newcommand{\Cov}{\operatorname{Cov}}
\newcommand{\supp}{\operatorname{supp}}
\DeclareMathOperator*{\argmin}{arg\,min}
\DeclareMathOperator*{\argmax}{arg\,max}
\newcommand{\entrymeta}[3]{{\sffamily\small\color{muted}\textbf{\texttt{#1}}\quad|\quad #2\par #3\par}}
\newcommand{\Problemref}[1]{\texttt{#1}}
\newcommand{\JELref}[2]{\texttt{#2} (JEL \texttt{#1})}
\begin{document}
\section*{Aid that builds lasting capacity}
\textbf{Problem ID:} \texttt{F35-575}\quad \textbf{JEL:} \texttt{F35}\par
% BEGIN ECONOMICS STATEMENT
\label{problem:OP-0375}
{\sffamily\small\color{muted}Original subject group: Development and poverty\par}
\label{definitions:problem:30007657}
\textbf{Audit classification:} Background; proposed extension. Full published version and metadata checked. Positive eleven-year results already exist; residual support continued, so these are not estimates after complete donor withdrawal.
\subsection*{Definitions and precise research target}
\textbf{Complete editorial problem.} Fix a set of aid-eligible communities and a program with a precisely declared withdrawal path beginning at date \(t_0\). Let \(p\) specify funding, local participation, implementer, public-service responsibilities, and withdrawal rules. Let \(K_{it}(p)\) measure locally maintained administrative or productive capacity and \(Y_{it}(p)\) household economic welfare. For a declared horizon \(T>t_0\), define \(\tau_K(T,p)=E[K_{iT}(p)-K_{iT}(0)]\) and \(\tau_Y(T,p)=E[Y_{iT}(p)-Y_{iT}(0)]\). Capacity must be measured using verified service delivery, staffing, maintenance, or collective-action outputs, rather than continued external spending. Estimate whether aid complements state and community investment or substitutes for it, allowing political incentives and nonbeneficiary spillovers. Randomized introduction identifies offer effects; withdrawal comparisons need credible timing variation and controls for concurrent shocks. Separate temporary assets from institutional changes by measuring outcomes after each specified support component ends. Determine which program features predict durable gains and can be transported to other delivery systems. A known long-run null or positive result in one setting is a partial answer. The present statement proposes a general identification and design extension; no exact primary passage establishing this complete question as open was verified in the supporting long-run aid study.

\textbf{Definitions and admissible scope.} Communities, rather than households alone, are the assignment units for capacity outcomes. Let \(K_{jt}\) be an independently measured vector of maintenance, staffing, tax collection, and service outputs, with a declared scalar index only when needed. Distinguish withdrawal of grants, withdrawal of facilitators, and withdrawal of all donor-paid personnel. The complete regime includes any residual support after \(t_0\). A post-withdrawal contrast requires both treatment and comparator to have specified support paths; calendar time since program launch is not equivalent to time since complete withdrawal. The research problem is identifying durable program effects and comparing delivery arrangements, not denying the positive long-run results already found in Sierra Leone. Let communities be indexed \(j\) with a fixed baseline household sample within each community. The community capacity target is \(\tau_K=E[K_{jT}(p)-K_{jT}(0)]\), componentwise for vector \(K\), and the household welfare target uses a declared household weighting law. The index \(i\) in the introductory shorthand is replaced by \(j\) for capacity units. Fix \(t_0<T<\infty\) and support paths for grants, facilitators, and personnel in both arms. For a scalar durability criterion choose a fixed \(v(K)\) and compare \(E[v(K_{jT}(p))-v(K_{jT}(0))]\); no ordering of a vector is assumed silently. Introduction and withdrawal assignment, follow-up selection, and community interference are explicit inputs.
\subsection*{Variants and assumption changes}
\begin{enumerate}
\item Observed long-run aid version (source-based scope): retain the actual support history of the Sierra Leone program and compare its treated and control communities after eleven years. This is an existing estimated contrast, not a newly open question.
\item Withdrawal design (editorial): randomize or credibly instrument grant and technical-assistance withdrawal dates separately. Measure locally funded maintenance and state service provision after each support component ends.
\end{enumerate}
\subsection*{References and source audit}
\textbf{Support and access.} Full published version and metadata checked. Positive eleven-year results already exist; residual support continued, so these are not estimates after complete donor withdrawal. \textbf{Locator:} Published 2023 paper, abstract and Section 3, pp. 1368--1369.
\begin{enumerate}\raggedright
\item\hypertarget{definitions:source:30007657:1}{} Katherine Casey; Rachel Glennerster; Edward Miguel; Maarten J. Voors. 2023. \emph{Long-Run Effects of Aid: Forecasts and Evidence from Sierra Leone}. The Economic Journal 133(652), 1348--1370; abstract and Section 3.
\url{https://emiguel.econ.berkeley.edu/wordpress/wp-content/uploads/2023/09/uead001.pdf}.
DOI: \url{https://doi.org/10.1093/ej/uead001}.
\end{enumerate}

\subsection*{Common conventions: Source mathematical conventions}

These conventions apply to each entry unless explicitly replaced.

\textbf{Models and identification.} A primitive model \(m\) specifies all probability laws, feasible choices, information, equilibrium and measurement rules used by its target. The admissible class \(\mathcal M\) is fixed independently of outcomes. If \(P_O(m)\) is its observable probability law and \(\tau(m)\) the target, its identified set at observable law \(P\) is
\[
 \mathcal I_\tau(P)=\{\tau(m):m\in\mathcal M,\ P_O(m)=P\}.
\]
Point identification means that this set is a singleton. An empty set signals incompatibility of data and assumptions. Coordinate bounds need not describe the joint set. Bounds are sharp only if every retained target is attainable or arbitrarily approachable by compatible models. Finite bounds require explicit restrictions on unobserved quantities; unrestricted confounding, hidden wealth, extrapolation or arbitrary equilibrium selection can yield uninformative sets.

\textbf{Causal interventions.} A potential outcome \(Y_i(p)\) is the outcome under a complete policy regime \(p\), including timing, financing, information and market spillovers. The target population and units are fixed before treatment. With interference, use the complete assignment vector or a justified exposure mapping; individual no-interference notation alone cannot identify an economy-wide intervention. A difference of mean potential outcomes does not require the cross-world joint distribution. Conditioning on post-treatment migration, survival, enrollment or employment changes the estimand and needs separate justification.

\textbf{Statistical statements.} An estimator, confidence set, forecast or test specifies a sampling law, model class and loss/coverage criterion. Uniform \((1-\alpha)\)-coverage means \(\inf_{m\in\mathcal M}\Pr_m\{\tau(m)\in C_n\}\geq1-\alpha\), or an explicitly asymptotic version. The whole parameter space trivially has coverage, so informativeness requires a stated width, risk or power objective. A forecasting improvement is relative to a named benchmark, information set and evaluation population.

\textbf{Equilibrium and optimization.} An equilibrium comprises feasible plans, optimality, beliefs where needed, budget constraints and all market/resource clearing. The primitives or a stated benchmark define each correspondence \(\mathcal E(p;m)\). Nonemptiness is assumed or proved, never inferred just from compact policy sets. A selected-equilibrium target uses a declared selection rule; a robust target quantifies over all permitted equilibria. Use suprema or infima unless attainment is established. An \(\varepsilon\)-optimal policy is feasible and within \(\varepsilon\) of the finite optimum value. Report an empty feasible set separately.

\textbf{Welfare and accounting.} Normative utility, interpersonal normalization, social weights, numeraire, discounting and the use of public receipts are primitives. Behavioral utility need not coincide with the welfare criterion. Household welfare includes ownership income and rents once; adding profits again to compensating variation generally double-counts them. A monetary equivalent is defined by an explicit transfer and price convention, with monotonicity and range conditions. Average effects, mechanism contributions, transfers and welfare losses are different targets. Infinite sums require convergence and dynamic feasible plans require terminal or no-Ponzi conditions.

\textbf{Source versions.} A working-paper date, revision date and journal publication year are distinguished. A DOI for a journal version is not automatically the DOI of an earlier manuscript. Locators refer to the stated version and printed pages where specified. A metadata-only check does not verify a claim about a full-text conclusion. Every entry's source subsection narrows the provenance claim accordingly.
% END ECONOMICS STATEMENT
\end{document}
